\chapter{Authentication Layers in zkVM Workflows}

\section{How the Current Encryption Method Was Chosen}

The final method was chosen in two stages. First, primitive benchmarks compared
AES, Salsa20, and LowMC in the same proving stack. Second, the chosen primitive
was integrated into an authenticated pipeline with explicit framing and
verifiable execution.

AES became the practical base because it performed better than LowMC in this
implementation context, and CTR mode offered a straightforward composition path
with HMAC in Encrypt-then-MAC style under NIST SP 800-38A, RFC 2104, and FIPS
198-1 \cite{nistsp80038a,rfc2104,nistfips1981}.
Chapter 6 records the cipher-selection decision; this chapter focuses on the
authentication layers required once AES-CTR is used as the confidentiality base.

\section{Threat-Driven Requirements}

The authentication design is driven by four concrete threats:

\begin{itemize}
\item an observer should not learn plaintext from public ciphertext;
\item an adversary should not be able to alter ciphertext, IV, or AAD without tag
verification failing;
\item a verifier should not accept outputs from a different guest program image;
\item a deployment should not confuse data authenticity with sender identity.
\end{itemize}

AES-CTR addresses the first threat under nonce-respecting use. HMAC over
\texttt{aad || iv || ciphertext} addresses the second under standard MAC-security
assumptions, provided tag verification is enforced before decryption. Receipt
verification against the expected image ID addresses the third. The fourth is
intentionally left outside the implemented cryptographic payload path and is
listed as a transport-layer limitation.

\section{Layer Separation}

This dissertation separates three layers:

\begin{itemize}
\item \textbf{Data authentication:} HMAC over framed data
(\texttt{aad || iv || ciphertext}) detects tampering in the ciphertext path.
\item \textbf{Computation correctness:} zkVM proof verification binds public
outputs to correct execution of a fixed guest image.
\item \textbf{Transport/sender identity:} external signatures or channel
authentication can bind packets to an operational identity.
\end{itemize}

This separation resolves the apparent chicken-and-egg issue. The proof does not
authenticate a network sender, and the MAC does not prove that the zkVM program
executed correctly. Instead, the MAC authenticates the framed data tuple, while
the receipt proves that the expected guest image produced or checked that tuple.
Both checks are therefore complementary rather than circular.

In short, a valid proof does not replace message authentication, and a valid MAC
does not prove correct execution of a fixed guest image. Both checks are needed
for the scoped threat model.

\section{Why Encrypt-then-MAC Runs in the VM}

Encrypt-then-MAC runs in the guest rather than as a host-side post-processing
step because the tag is part of the statement being proved. The verifier should
be able to interpret a receipt as evidence that the fixed guest image encrypted
the plaintext, authenticated \texttt{aad || iv || ciphertext}, and committed the
resulting public transcript. This gives a stronger receipt-binding claim than
proving encryption first and letting the host attach a tag afterwards.

If HMAC were computed only outside the VM, the receipt would no longer bind the
tag to the guest execution trace. A verifier could learn that some guest program
produced a ciphertext and that some host later produced a valid MAC, but not that
the guest computed the authenticated transcript under the stated framing rule.
That weaker claim is not the one used by the security game in Chapter 10, where
the public tuple contains \((\mathrm{aad}, n, c, t, y, \pi)\) and the proved
relation includes
\(t = \mathrm{Trunc}_{192}(\mathrm{HMAC}_{k_M}(\mathrm{aad} \parallel n \parallel c))\).

This placement also clarifies where optional transport authentication belongs. A
deployment may still sign or MAC the final receipt and packet outside the VM to
identify the sender or protect a network channel. That outer authentication is
not a substitute for the in-guest MAC; it answers who sent the packet, while the
in-guest Encrypt-then-MAC path answers whether the ciphertext/tag pair belongs to
the proved computation.

Input preparation, proving orchestration, receipt verification, and benchmark
artifact collection run in the host. Keeping these responsibilities in host code
matches the RISC Zero execution model and keeps deployment concerns separate
from deterministic guest logic \cite{risczeroDocs}.

\section{Alternatives Considered}

\paragraph{AES-GCM.}
AES-GCM combines encryption and authentication in one mode and is widely used,
but this project favored an explicit EtM split because it produced cleaner
component-level overhead attribution and a simpler game-based argument structure
for the dissertation: AES-CTR, HMAC, and the RISC Zero proof each correspond to
one hop in the Chapter 10 IND-CPA game
\cite{nistsp80038d,kasper2009aesgcm,bellare2000authenticated}.

\paragraph{ChaCha20-Poly1305.}
ChaCha20-Poly1305 is a strong software baseline with broad protocol adoption,
but it was not selected here because the repository and benchmark pipeline were
already centered on AES-derived targets and because introducing a new primitive
family would dilute depth in the final report campaign \cite{rfc8439}.

\paragraph{OCB mode.}
OCB provides efficient one-pass authenticated encryption, but this project
preferred conservative, widely standardized components under explicit composition
reasoning for easier implementation audit and chapter-level explanation
\cite{rfc7253,bellare2000authenticated}.

\begin{table}[t]
\centering
\scriptsize
\begin{tabular}{p{0.16\linewidth}p{0.21\linewidth}p{0.21\linewidth}p{0.21\linewidth}}
\hline
Alternative & Security fit & Evaluation fit & Reason not selected \\
\hline
AES-GCM & Strong AEAD when nonce discipline is correct & Would require
GHASH-specific zkVM measurement & Larger integrated surface and less direct
CTR-vs-auth overhead attribution \\
ChaCha20-Poly1305 & Strong protocol baseline with software-friendly operations &
Would introduce a new primitive family late in the project & Reduces depth of
AES/LowMC comparison and final artifact traceability \\
OCB & Efficient one-pass AEAD with formal support & Not already implemented or
benchmarked in this repository & Additional implementation and licensing/history
discussion not central to research question \\
AES-CTR + HMAC & Standard components with clear EtM reasoning & Directly
supports separate CTR and CTR+HMAC targets & Selected despite two-pass overhead \\
\hline
\end{tabular}
\caption{Authentication-mode alternatives compared against project requirements.}
\label{tab:auth-alternatives}
\end{table}

\section{Current Implementation Scope}

The implemented path covers the first two layers (HMAC plus zkVM proof
verification). Additional transport-level proof authentication is outside the
current code path.

Therefore, report claims are scoped to confidentiality and integrity of the
cryptographic computation transcript under the fixed-program verification model,
not to sender identity guarantees in a network deployment. They also assume
nonce/IV uniqueness for AES-CTR and do not claim safety if the same counter stream
is reused under the same encryption key.
